\chapter{Dry Mouth Syndrome}
\index{Dry Mouth Syndrome}

\section*{Definition and Overview}
Dry mouth, also known as xerostomia, is the unpleasant sensation of reduced saliva in the mouth. Saliva is essential for oral health: it lubricates food, protects the teeth against decay, neutralises acids, and keeps the mouth comfortable and free of infection. Dry mouth is not a disease in itself but a symptom of an underlying cause, most commonly medications, and it is very common, particularly in older people. This chapter explains the causes of dry mouth, the problems it can cause, and the many effective ways to manage it.

\section*{Epidemiology}
Dry mouth is extremely common, affecting a significant proportion of the general population and a much higher proportion of older people. It is estimated that about one in five older adults experiences persistent dry mouth, and the figure is higher still among people in residential care. The condition is far more prevalent in people taking multiple medicines, which is most older adults, and in those with autoimmune diseases, diabetes, or a history of head and neck radiotherapy. Dry mouth is often accepted as an unavoidable part of ageing or medication, and many people do not mention it to their doctor, despite the significant effects it can have on comfort, eating, and dental health.

\section*{Aetiology and Pathophysiology}
Saliva is produced by the parotid, submandibular, and sublingual glands, along with many small glands throughout the mouth. Dry mouth results when saliva production falls below what is needed, or when saliva is lost through breathing or evaporation.
\begin{itemize}
    \item \textbf{Medications:} the most common cause; hundreds of medicines reduce saliva, including antihistamines, decongestants, antidepressants, antipsychotics, diuretics, and medicines for overactive bladder and high blood pressure
    \item \textbf{Dehydration:} inadequate fluid intake, fever, vomiting, diarrhoea, or excessive sweating
    \item \textbf{Ageing:} the glands become less productive, and older people take more medications, although dry mouth is not an inevitable part of ageing
    \item \textbf{Medical conditions:} Sjögren's syndrome, diabetes, HIV, and autoimmune disorders can damage the salivary glands
    \item \textbf{Radiotherapy:} radiation to the head and neck can permanently reduce saliva production
    \item \textbf{Chemotherapy:} may temporarily affect saliva production and composition
    \item \textbf{Lifestyle factors:} smoking, alcohol, and mouth breathing, especially at night
    \item \textbf{Nerve damage:} injury to the head or neck can affect the nerves supplying the glands
\end{itemize}

\section*{Clinical Features}
The symptoms of dry mouth range from a mild, manageable dryness to severe discomfort.
\begin{itemize}
    \item A persistent feeling of stickiness or dryness in the mouth
    \item Difficulty chewing, swallowing, and speaking
    \item Changes in taste, with food tasting bland or unusual
    \item A sore or burning tongue
    \item Cracked lips and sores at the corners of the mouth
    \item Bad breath
    \item A dry, scratchy throat and hoarse voice
    \item Difficulty wearing dentures
    \item Increased thirst
\end{itemize}

\section*{Differential Diagnosis}
Because dry mouth is a symptom, the priority is identifying the cause.
\begin{itemize}
    \item Medication-induced dry mouth, the most common explanation
    \item Dehydration and inadequate fluid intake
    \item Sjögren's syndrome and other autoimmune conditions
    \item Diabetes, which can cause thirst and mouth dryness
    \item Mouth breathing, particularly overnight with snoring or sleep apnoea
    \item Anxiety and stress, which can reduce saliva flow
\end{itemize}

\section*{When to Seek Medical Care}
See a GP or dentist if dry mouth is persistent, affecting eating or speaking, or accompanied by dental problems. A doctor can review medications, check for underlying conditions, and arrange referral to a specialist where needed. Dentists can help protect the teeth from the effects of reduced saliva.

\textbf{Contact your local emergency services for:}
\begin{itemize}
    \item Difficulty breathing or severe swelling of the mouth or throat
    \item Anaphylaxis symptoms, including throat tightness, wheezing, and dizziness
    \item Sudden inability to swallow or speak
\end{itemize}

\section*{Diagnosis and Investigations}
Diagnosis of dry mouth is based on the history and examination, with tests used to identify underlying causes.
\begin{itemize}
    \item \textbf{History:} medication list, fluid intake, symptoms, and any underlying conditions
    \item \textbf{Examination:} inspection of the mouth for dryness, dental decay, and signs of infection; checking the flow of saliva from the gland openings
    \item \textbf{Saliva measurement:} simple tests to estimate saliva production
    \item \textbf{Blood tests:} to check for diabetes, Sjögren's syndrome, and other autoimmune or inflammatory conditions
    \item \textbf{Salivary gland imaging or biopsy:} in selected cases where gland damage is suspected
    \item \textbf{Medication review:} identifying and, where possible, adjusting medicines that cause dryness
\end{itemize}

\section*{Management}
Management addresses the cause where possible and relieves symptoms in the meantime.
\begin{itemize}
    \item \textbf{Medication adjustment:} a doctor may change the dose, timing, or choice of medicines contributing to dry mouth
    \item \textbf{Frequent sips of water:} the simplest and most effective measure for most people
    \item \textbf{Saliva substitutes:} sprays, gels, and lozenges that mimic saliva and moisten the mouth
    \item \textbf{Saliva stimulants:} sugar-free chewing gum or sweets, and prescription medications that stimulate saliva production
    \item \textbf{Good oral hygiene:} regular brushing, flossing, and fluoride use to protect the teeth against decay
    \item \textbf{Dietary adjustments:} soft, moist foods and avoidance of salty, spicy, or sugary items
    \item \textbf{Humidifying the air:} using a humidifier, especially in the bedroom
    \item \textbf{Avoiding irritants:} reducing alcohol, caffeine, tobacco, and mouthwashes containing alcohol
    \item \textbf{Dental care:} regular check-ups and fluoride treatments to prevent decay
\end{itemize}

\section*{Complications}
\begin{itemize}
    \item Dental caries, because saliva's protective and buffering actions are reduced
    \item Gum disease and oral infections, including candidiasis (thrush)
    \item Difficulty chewing, swallowing, and speaking, affecting nutrition and communication
    \item Cracked lips, mouth ulcers, and a burning sensation
    \item Reduced enjoyment of food and social eating
    \item Denture problems, including soreness and poor fit
\end{itemize}

\section*{Prognosis and Outcomes}
The outlook for dry mouth depends on the cause. Dryness caused by medications, dehydration, or lifestyle factors usually improves substantially when the cause is addressed. Dryness from radiotherapy or autoimmune disease may be permanent and requires ongoing management, but good symptom control is achievable for most people. Regular dental care is essential, because the risk of decay and infection continues while saliva is reduced.

\section*{Prevention and Self-Care}
\begin{itemize}
    \item Drink water regularly throughout the day, and keep a bottle within reach
    \item Chew sugar-free gum or suck sugar-free sweets to stimulate saliva
    \item Use a saliva substitute spray or gel if dryness persists
    \item Avoid alcohol, tobacco, and caffeine, which worsen dryness
    \item Keep lips moisturised with balm
    \item Maintain meticulous oral hygiene and see the dentist regularly
    \item Discuss any dry mouth with your doctor when starting new medications
\end{itemize}

\section*{Sources and Further Reading}
The following reputable sources provide further information for healthcare students and patients seeking reliable, current advice about dry mouth.
\begin{itemize}
    \item NHS. \textit{Dry mouth (xerostomia): causes and treatment.} https://www.nhs.uk/conditions/dry-mouth/
    \item Mayo Clinic. \textit{Dry mouth: symptoms, causes, and treatment.} https://www.mayoclinic.org/diseases-conditions/dry-mouth/
    \item MedlinePlus. \textit{Dry mouth.} https://medlineplus.gov/drymouth.html
    \item Australian Dental Association. \textit{Oral health information on dry mouth.} https://www.ada.org.au/
    \item National Institute of Dental and Craniofacial Research. \textit{Xerostomia patient information.} https://www.nidcr.nih.gov/
    \item MSD Manual. \textit{Dry mouth: professional overview.} https://www.msdmanuals.com/
\end{itemize}
